\documentclass[conference]{IEEEtran}
\usepackage{cite}
\usepackage{amsmath,amssymb,amsfonts}
\usepackage{algorithmic}
\usepackage{graphicx}
\usepackage{textcomp}
\usepackage{xcolor}
\usepackage{booktabs}
\usepackage{array}
\usepackage{url}
\usepackage{hyperref}
\usepackage{microtype}

\def\BibTeX{{\rm B\kern-.05em{\sc i\kern-.025em b}\kern-.08em
    T\kern-.1667em\lower.7ex\hbox{E}\kern-.125emX}}

\begin{document}

\title{Persistent Repository Memory for Autonomous Software Engineering Agents: Architectural Foundations, Empirical Retrieval, and Downstream Resolution Limits}

\author{
\IEEEauthorblockN{
Abhinav Pant\IEEEauthorrefmark{1},
A Anjana Nair\IEEEauthorrefmark{2},
Abhiram Anil\IEEEauthorrefmark{3},
and Kalayraman P\IEEEauthorrefmark{4}
}
\IEEEauthorblockA{
Department of Computer Science and Engineering
}
}

\maketitle

\begin{abstract}
Improvements in repository retrieval do not automatically translate into improved coding performance. We study this gap through a \textbf{Session-Persistent Repository Memory Architecture} for autonomous software engineering (SWE) agents, which are normally stateless: each issue runs in an isolated container that discards structural indexes, debugging trajectories, and execution feedback. Our architecture combines (1) Tree-sitter symbol graphs with subtoken decomposition; (2) Reciprocal Rank Fusion (RRF, $k=60$) of sparse lexical retrieval (sublinear TF-IDF) and AST symbol matching, with no GPU embedding overhead; (3) SHA-256 content-hash caching that skips redundant parsing ($10.13\times$ speedup); (4) scoped episodic outcome logging under a zero-leakage constraint; and (5) SHA-256 staleness invalidation that attenuates outdated episodes by 80\% ($\lambda=0.20$).

We evaluate on \textbf{SWE-bench Verified Mini} ($N=50$: 25 \texttt{django/django}, 25 \texttt{sphinx-doc/sphinx}), checking out each task's exact \texttt{base\_commit}. Hybrid RRF improves fault localization over snapshot lexical retrieval on Django (MRR 0.6027 vs.\ 0.3667, +64.4\%; paired Wilcoxon $p=0.0071$ two-sided; Top-5 hit rate 80.0\%) and on the combined benchmark (MRR 0.5213 vs.\ 0.4163, +25.2\%; Top-5 hit rate 70.0\%). The gain is repository-dependent: on Sphinx, pure AST retrieval performs poorly, and under a single-snapshot protocol Hybrid RRF has a lower MRR than lexical retrieval (0.2447 vs.\ 0.3033) while improving Top-5 hit rate (44.0\% vs.\ 40.0\%). Scoped episodic memory reduced exploratory searches in two recurrent Django subsystem clusters (15$\to$5 and 5$\to$0 searches), but \textbf{we do not demonstrate an improvement in task resolution}: in the $N=7$ end-to-end runs the proposed system resolved 2/7 tasks versus 3/7 for the stateless baseline, even when the target file was ranked first. Unscoped episodic memory causes negative transfer (MRR falls from 0.5213 to 0.1860), and dense structural context can induce over-deliberation (one observed failure). In a four-scenario drift study, SHA-256 tracking invalidated obsolete definitions; the $\lambda$ sweep shows only that hard eviction ($\lambda=0$) is harmful, not that $\lambda=0.20$ is optimal.
\end{abstract}

\begin{IEEEkeywords}
Autonomous Coding Agents, Software Engineering, Persistent Memory, AST Symbol Graphs, Reciprocal Rank Fusion, SWE-bench Verified.
\end{IEEEkeywords}

\section{Introduction}
Autonomous coding agents (e.g., SWE-agent~\cite{yang2024sweagent}, Devin, Agentless~\cite{xia2024agentless}, and mini-swe-agent~\cite{lieret2025minisweagent}) have demonstrated remarkable capabilities on software engineering benchmarks such as SWE-bench~\cite{jimenez2024swebench}. Nevertheless, standard agent scaffolding remains strictly \textbf{stateless}: every issue starts from an empty environment, forcing the agent to discover repository structure, file layouts, and subsystem conventions entirely from scratch.

In large codebases containing hundreds of thousands of lines of code (such as the Django framework), statelessness incurs severe operational costs:
\begin{enumerate}
    \item \textbf{Unguided Exploration}: Without structural context, agents spend dozens of turns executing exploratory bash commands (\texttt{find}, \texttt{grep}) to locate relevant files, frequently exhausting turn budgets before reaching code synthesis.
    \item \textbf{Episodic Amnesia}: Insights gained during previous task debugging (such as module conventions or test script locations) are discarded upon container termination.
    \item \textbf{Repository Drift Vulnerability}: In persistent workflows where code evolves across commits, agents caching historical trajectories risk relying on obsolete symbol signatures unless cryptographic staleness mechanisms are in place.
\end{enumerate}

Surveys of software-engineering agents and LLM-based multi-agent systems show a decisive shift toward architectures combining perception, memory, planning, tool use, and action~\cite{wang2025agents, he2025multiagent}. However, existing memory mechanisms are largely conversational (interaction-oriented) or artifact-based (persisting bug report summaries), rather than being grounded in the concrete syntax and evolving history of the source-code repository itself.

This paper presents an implementation and empirical evaluation of a \textbf{Session-Persistent Repository Memory Architecture}. Our central finding is that improved retrieval is not equivalent to improved coding performance:
\begin{itemize}
    \item \textbf{Localization improves, with repository-dependent gains}: Fusing Tree-sitter AST symbol graphs with sparse lexical retrieval via RRF ($k=60$) improves fault localization on Django ($N=25$; MRR 0.6027 vs.\ 0.3667 under task-specific checkouts, two-sided Wilcoxon $p=0.0071$) and on the combined 50-task benchmark. On Sphinx the evidence is mixed (Section~\ref{sec:retrieval}).
    \item \textbf{Scoped episodic transfer reduces exploration in the clusters we tested}: In two recurrent Django subsystem pairs (backend ORM and frontend widgets, $N=2$ task pairs), persistent memory reduced exploratory searches (15$\to$5 and 5$\to$0) and, on \texttt{django-12276}, shortened a verified resolution by 78 seconds (a single task).
    \item \textbf{The Localization--Resolution Disconnect}: Accurate Rank-1 localization did not raise resolution rates. Across $N=7$ end-to-end tasks the proposed system resolved 2/7 versus 3/7 for the stateless baseline. We therefore make no claim of task-success improvement.
    \item \textbf{Staleness verification safeguards persistent state}: SHA-256 verification detects on-disk modifications and down-weights outdated context by 80\%, addressing a vulnerability of naive string-logging memory (evaluated on four controlled drift scenarios).
\end{itemize}

\section{Literature Foundation \& Taxonomy}
Our work builds on foundational peer-reviewed studies published across empirical software engineering, artificial intelligence, and automated software engineering venues:

\subsection{History-Aware Repository Context}
\textbf{HAFix}~\cite{shi2026hafix} directly addresses the limitation that LLM-based bug-fixing approaches rely primarily on the current project snapshot while overlooking historical repository information. HAFix introduces history-augmented heuristics and aggregates historical commit signals to improve bug fixing. Its significance is conceptual as much as empirical: repository history is not merely archival metadata; selected historical evidence can become actionable task context. However, HAFix supplies historical heuristics for isolated single-task bug-fixing formulations, rather than maintaining a continuously evolving, multi-task repository memory representation.

\subsection{Repository-Level Retrieval \& Multi-Agent Refactoring}
\textbf{RepoAI}~\cite{chondamrongkul2026repoai} demonstrates repository-level code understanding and refactoring through specialized agents combined with retrieval-augmented generation. Complementary methods such as RepoCoder~\cite{zhang2023repocoder} explore iterative retrieval for code completion. Wang et al.~\cite{wang2026agr} introduce \textbf{Agent-Governed Retrieval-Augmented Generation (AGR)}, which retrieves evidence from documentation, historical reports, and web sources, using an agentic gate to judge evidence sufficiency. However, RepoAI is framed around refactoring pipelines, and AGR persists report-level artifacts rather than maintaining incrementally updated, code-structural representations of source entities and interaction outcomes across sequential coding tasks.

\subsection{Software-Engineering Agents, AI-IDEs, and Memory}
The comprehensive survey by \textbf{Wang et al.}~\cite{wang2025agents} organizes software-engineering agents around perception, memory, and action, establishing memory as a first-class architectural component. \textbf{He et al.}~\cite{he2025multiagent} extend this perspective to LLM-based multi-agent systems, highlighting collaboration, specialized roles, and open challenges. 

\textbf{Contreras et al.}~\cite{contreras2026aiides} architect open and trustworthy AI-IDEs where developer-assistant interactions are persistently traced alongside their impact on code, accompanied by quality-assurance validation loops (compilation, static analysis, testing). While relevant to interaction memory, their traceability layer is not an adaptive repository-memory retrieval mechanism that dynamically surfaces structural code evidence for future coding tasks.

From the foundational AI perspective, \textbf{The AI Hippocampus}~\cite{jia2025hippocampus} formalizes memory in LLM and multimodal systems into implicit, explicit, and agentic paradigms. Explicit memory includes queryable external representations (text, vectors, graphs), while agentic memory emphasizes persistent, temporally extended structures. In a repository setting, structural code graphs serve as explicit memory, while task execution histories and test outcomes provide an agentic memory layer.

\subsection{Program Repair \& Execution-Grounded Validation}
Empirical work on LLM-based program repair demonstrates that prompt context design materially dictates repair effectiveness~\cite{cao2025programrepair}. Furthermore, \textbf{Jin \& Chen}~\cite{jin2026codereview} demonstrate that LLM-based requirement-conformance judgments exhibit systematic overcorrection errors and hallucinations. They establish that agent judgments must be grounded in executable tests and external verification. This justifies associating stored interaction outcomes with SWE-bench harness execution evidence rather than relying on model self-assessment.

\section{Capability Matrix \& Research Gaps}

\subsection{Structured Capability Comparison}
Table~\ref{tab:comparison} compares how prior state-of-the-art architectures handle key dimensions of repository memory compared to our proposed architecture. The comparison is qualitative, based on each paper's described scope; we do not re-implement these systems as baselines.

\begin{table*}[t]
\caption{Structured Capability Comparison Across Literature and Proposed Architecture}
\label{tab:comparison}
\centering
\setlength{\tabcolsep}{4.5pt}
\resizebox{\textwidth}{!}{%
\begin{tabular}{lcccccc}
\toprule
\textbf{System / Paper} & \textbf{Persistent Memory} & \textbf{Code Dependencies} & \textbf{Repo Evolution} & \textbf{Stale-Info Handling} & \textbf{Multi-Turn} & \textbf{Retrieval Efficiency} \\
\midrule
HAFix~\cite{shi2026hafix} & No & No & Partial (Git) & No & No & No \\
RepoAI~\cite{chondamrongkul2026repoai} & No & Partial & No & No & Partial & No \\
Wang et al.~\cite{wang2025agents} & Partial & No & No & No & Partial & No \\
He et al.~\cite{he2025multiagent} & Partial & No & No & No & Partial & No \\
The AI Hippocampus~\cite{jia2025hippocampus} & Partial & No & No & No & Partial & No \\
AGR~\cite{wang2026agr} & Partial (Reports) & No & Partial & Partial & Partial & No \\
AI-IDEs~\cite{contreras2026aiides} & Yes (Traces) & Partial & Partial & No & Yes & No \\
Jin \& Chen~\cite{jin2026codereview} & No & No & No & No & No & No \\
Cao et al.~\cite{cao2025programrepair} & No & Partial & No & No & Partial & No \\
\midrule
\textbf{Proposed System (This Work)} & \textbf{Yes (Episodic)} & \textbf{Yes (AST Graphs)} & \textbf{Yes (Incremental)} & \textbf{Yes (SHA-256)} & \textbf{Yes (Docker)} & \textbf{Yes (RRF + Cache)} \\
\bottomrule
\end{tabular}%
}
\end{table*}

\subsection{Formal Research Gaps}
From this landscape, we synthesize five core research gaps:
\begin{itemize}
    \item \textbf{Gap 1: Historical repository knowledge is not yet a persistent memory substrate}: Prior systems supply historical heuristics for isolated single tasks rather than maintaining an evolving structural representation.
    \item \textbf{Gap 2: Repository RAG and persistent memory remain unintegrated}: Existing RAG systems operate at task-time or persist textual report summaries rather than incrementally updated structural representations of code entities.
    \item \textbf{Gap 3: Unified representation of code structure, traceability, and memory}: Dependencies, call graphs, edits, retrieved context, and test outcomes must coexist in an incrementally maintainable data structure.
    \item \textbf{Gap 4: Long-horizon sequential evaluation is underdeveloped}: Agent benchmarks predominantly evaluate isolated tasks. Measuring the cumulative effect of stored memory requires chronological evaluation.
    \item \textbf{Gap 5: Retrieval must be adaptive and drift-resilient}: When source files evolve, naive memory risks serving outdated symbol signatures and invalid historical patches.
\end{itemize}

\section{System Architecture \& Mathematical Formulation}

\subsection{Tree-sitter AST Symbol Extraction \& Subtoken Decomposition}
Given a Python source file $f \in \mathcal{F}$, we parse its concrete syntax tree using Tree-sitter, building on lightweight repository mapping paradigms~\cite{gauthier2023aider} without the resource overhead of heavy dense pre-trained encoders~\cite{feng2020codebert, guo2021graphcodebert}. We extract top-level definitions:
\begin{equation}
\mathcal{S}(f) = \left( \mathcal{C}_f, \mathcal{M}_f, \mathcal{K}_f \right)
\end{equation}
where $\mathcal{C}_f$ denotes declared class identifiers, $\mathcal{M}_f$ denotes function/method definitions, and $\mathcal{K}_f$ denotes outgoing call invocations.

To match natural-language descriptions that reference partial or composite names, we apply subtoken decomposition supporting CamelCase, snake\_case, and uppercase acronyms:
\begin{equation}
\text{Subtokens}(w) = \{ \text{lower}(s) \mid s \in \text{RegexFindall}(w, \text{pattern}), |s| \ge 3 \}
\end{equation}
using regex pattern \texttt{"[A-Z]+(?=[A-Z][a-z0-9]|$\backslash$b)|[A-Z][a-z0-9]*|[a-z0-9]+"}.

\subsection{Hybrid Reciprocal Rank Fusion (RRF)}
Let $\mathcal{R}_{\text{lex}}(f)$ denote the rank of file $f$ under sparse lexical retrieval (sublinear TF-IDF cosine similarity), and $\mathcal{R}_{\text{ast}}(f)$ denote the rank of file $f$ under AST symbol matching. The unified fusion score is defined via Reciprocal Rank Fusion~\cite{cormack2009rrf}:
\begin{equation}
\text{RRF}(f) = \frac{1}{k + \mathcal{R}_{\text{lex}}(f)} + \frac{1}{k + \mathcal{R}_{\text{ast}}(f)}, \quad k = 60
\end{equation}
Files are sorted in descending order of $\text{RRF}(f)$, and top $K=5$ files are injected into the agent prompt along with signatures.

\subsection{Incremental AST Caching}
We maintain an in-memory symbol cache:
\begin{equation}
\text{Cache}: \text{rel\_path} \to (\text{SHA256}, \mathcal{S}(f))
\end{equation}
On task $t$, the system computes SHA-256 hashes for repository source files ($O(M)$ disk I/O, empirically $\approx 243\text{ ms}$ for 2,570 files in Django). If $\text{SHA256}_t(f) = \text{SHA256}_{t-1}(f)$, syntax tree parsing is bypassed ($O(1)$ cache hit).

\subsection{Cryptographic Staleness Invalidation}
When task $t$ completes, the system records:
\begin{equation}
e_t = \left( \text{ID}_t, \mathcal{P}_t, \mathcal{F}^{\text{edited}}_t, \mathcal{F}^{\text{retrieved}}_t, y_t, \mathcal{H}_t \right)
\end{equation}
where $y_t \in \{0, 1\}$ denotes harness resolution and $\mathcal{H}_t = \{ f \to \text{SHA256}(f) \}$ stores file hashes. On subsequent task $t' > t$, if $\exists f \in \mathcal{H}_t \text{ s.t. } \text{SHA256}_{t'}(f) \neq \mathcal{H}_t(f)$, the episode is marked stale. Base weight $w_0(e_t) \in \{+1.0, -0.5\}$ is attenuated via $\lambda = 0.20$:
\begin{equation}
w(e_t) = \begin{cases} w_0(e_t) & \text{if } \neg\text{Stale}(e_t) \\ \lambda \cdot w_0(e_t) & \text{if } \text{Stale}(e_t) \end{cases}
\end{equation}
Note that invalidation is file-granular: any change to a file in $\mathcal{H}_t$ marks the whole episode stale, even if the edit is unrelated to the symbols the episode concerns.

\section{Experimental Setup}
We evaluate on \textbf{SWE-bench Verified Mini}~\cite{jimenez2024swebench}: 25 tasks from \texttt{django/django} and 25 from \texttt{sphinx-doc/sphinx} ($N=50$). Retrieval is evaluated under two protocols:
\begin{itemize}
    \item \textbf{Task-specific checkout (primary)}: For every task, repository HEAD is pinned to that issue's historical \texttt{base\_commit} before indexing or retrieval, guaranteeing no future information leakage.
    \item \textbf{Single snapshot (supplementary)}: A shared host checkout pinned to commit \texttt{b1d6b35e} (Django 3.0-dev, corresponding to \texttt{django-11790}). This protocol is used for the component ablations and the episodic-transfer experiments.
\end{itemize}
Zero ground-truth patch leakage was confirmed across all tasks. End-to-end runs used \texttt{mini-swe-agent 2.4.6}~\cite{lieret2025minisweagent}, \texttt{swebench 4.1.0}, and \texttt{gemini-3.5-flash-lite} (temperature = 0.0) with an explicit 50-turn pacing budget prompt.

\subsection{Research Questions}
Our empirical study investigates five core research questions:
\begin{itemize}
    \item \textbf{RQ1 (Fault Localization)}: Does fusing Tree-sitter AST symbol graphs with sparse lexical retrieval via Reciprocal Rank Fusion ($k=60$) improve fault localization accuracy over baseline retrieval across repository tasks?
    \item \textbf{RQ2 (Multi-Turn Resolution vs. Localization)}: Does placing the target file at Rank 1 translate to higher end-to-end task resolution rates within a multi-turn SWE agent loop?
    \item \textbf{RQ3 (Incremental AST Caching Amortization)}: Does content-addressed AST symbol caching amortize parsing overhead across consecutive tasks without disk thrashing?
    \item \textbf{RQ4 (Cryptographic Staleness Tracking)}: Can SHA-256 change detection prevent agents from serving outdated memory trajectories when repository source files evolve?
    \item \textbf{RQ5 (Cross-Task Episodic Transfer)}: Does persistent episodic memory transfer actionable context and reduce exploratory tool overhead across sequential tasks within recurrent subsystems?
\end{itemize}

\section{Quantitative Results}
\label{sec:retrieval}

\subsection{Fault Localization under Task-Specific Checkouts ($N=50$, Answering RQ1)}
Table~\ref{tab:chrono} reports retrieval quality when each task is evaluated at its exact \texttt{base\_commit}. On Django, Static Hybrid RRF reaches an MRR of \textbf{0.6027}, a +64.4\% relative gain over snapshot lexical (BM25/TF-IDF) retrieval at 0.3667 (paired Wilcoxon $W=13.0$, two-sided $p=0.0071$, directional $p=0.0036$). Combined across all 50 tasks, Hybrid RRF reaches MRR 0.5213 versus 0.4163 for lexical retrieval (+25.2\%). Incremental AST caching re-parsed an average of only 38 files between successive checkouts and served over 750 files from cache.

\begin{table}[!t]
\caption{Fault Localization with Task-Specific Checkouts (Static Hybrid RRF)}
\label{tab:chrono}
\centering
\setlength{\tabcolsep}{3.5pt}
\resizebox{\columnwidth}{!}{%
\begin{tabular}{lcccc}
\toprule
\textbf{Repository} & \textbf{Hit@1} & \textbf{Hit@5} & \textbf{MRR (Hybrid)} & \textbf{MRR (Lexical)} \\
\midrule
django/django ($N=25$) & 48.0\% (12/25) & 80.0\% (20/25) & 0.6027 & 0.3667 \\
sphinx-doc/sphinx ($N=25$) & 36.0\% (9/25) & 60.0\% (15/25) & --- & --- \\
\textbf{Combined ($N=50$)} & \textbf{42.0\% (21/50)} & \textbf{70.0\% (35/50)} & \textbf{0.5213} & \textbf{0.4163} \\
\bottomrule
\end{tabular}%
}
\end{table}

\subsection{Single-Snapshot Multi-Repository Benchmark ($N=50$, Answering RQ1)}
To compare all retrieval variants under identical conditions, we also evaluated the 50 tasks against the single pinned snapshot. Table~\ref{tab:retrieval_multirepo} reports HitRate@K and MRR.

\begin{table}[!t]
\caption{Multi-Repository Fault Localization, Single-Snapshot Protocol ($N=50$ Tasks)}
\label{tab:retrieval_multirepo}
\centering
\setlength{\tabcolsep}{3.5pt}
\resizebox{\columnwidth}{!}{%
\begin{tabular}{lcccc}
\toprule
\textbf{Repository \& Condition} & \textbf{Hit@1} & \textbf{Hit@5} & \textbf{MRR} & \textbf{Hits@5} \\
\midrule
\multicolumn{5}{l}{\textit{django/django ($N=25$ Tasks)}} \\
B (Snapshot TF-IDF) & 0.2000 & 0.6000 & 0.3667 & 15 / 25 \\
C (Naive Combo) & 0.2000 & 0.6000 & 0.3667 & 15 / 25 \\
D (Pure AST Ablation) & 0.3200 & 0.6800 & 0.4447 & 17 / 25 \\
\textbf{D (Static Hybrid RRF)} & \textbf{0.4400} & \textbf{0.8000} & \textbf{0.5773} & \textbf{20 / 25} \\
\midrule
\multicolumn{5}{l}{\textit{sphinx-doc/sphinx ($N=25$ Tasks)}} \\
B (Snapshot TF-IDF) & 0.2400 & 0.4000 & 0.3033 & 10 / 25 \\
C (Naive Combo) & 0.2400 & 0.4000 & 0.3033 & 10 / 25 \\
D (Pure AST Ablation) & 0.0000 & 0.1600 & 0.0600 & 4 / 25 \\
D (Static Hybrid RRF) & 0.1600 & 0.4400 & 0.2447 & 11 / 25 \\
\midrule
\multicolumn{5}{l}{\textit{Combined Multi-Repository Benchmark ($N=50$ Tasks)}} \\
B (Snapshot TF-IDF) & 0.2200 & 0.5000 & 0.3350 & 25 / 50 \\
C (Naive Combo) & 0.2200 & 0.5000 & 0.3350 & 25 / 50 \\
D (Pure AST Ablation) & 0.1600 & 0.4200 & 0.2523 & 21 / 50 \\
\textbf{D (Static Hybrid RRF)} & \textbf{0.3000} & \textbf{0.6200} & \textbf{0.4110} & \textbf{31 / 50} \\
\bottomrule
\end{tabular}%
}
\end{table}

\textbf{Statistical Significance (Django, snapshot)}: A paired Wilcoxon signed-rank test on MRR confirms a significant +57.4\% improvement ($W = 13.5$, two-sided $p = 0.0077$, directional $p = 0.0039$, $n=15$ non-zero pairs). Directional McNemar tests confirm improvements on HitRate@1 ($p = 0.035$) and HitRate@5 ($p = 0.031$; Wilson 95\% CI [0.609, 0.911] vs.\ [0.407, 0.766]). Combined across $N=50$, Hybrid RRF improves MRR to 0.4110 (+22.7\%) and Top-5 hit rate to 62.0\% (31/50), versus 50.0\% (25/50) for lexical and 42.0\% (21/50) for pure AST. We note that significance tests are reported for Django only; no test is reported for the Sphinx or combined results.

\textbf{Architectural Divergence and the Sphinx Result}: In object-oriented systems with explicit symbol inheritance (Django), AST symbols align with issue descriptions, and pure AST reaches 0.4447 MRR. In compiler/DSL architectures (Sphinx), problem statements frequently reference reStructuredText directive syntax rather than Python identifiers; pure AST drops to 0.0600 MRR while lexical retrieval reaches 0.3033. Under the single-snapshot protocol, Hybrid RRF on Sphinx improves Top-5 hit rate over lexical (44.0\% vs.\ 40.0\%) but has \emph{lower} Hit@1 (16.0\% vs.\ 24.0\%) and MRR (0.2447 vs.\ 0.3033). Fusion therefore helps recall in the top five but does not uniformly dominate lexical retrieval; the benefit of AST signals depends on how closely issue text references code identifiers.

\subsection{Architectural Component-Wise Ablations (Answering RQ1 \& RQ4)}
Table~\ref{tab:ablations} reports isolated component contributions across lexical, structural, static fusion, and unfiltered episodic mechanisms (single-snapshot protocol; values for variants (1) and (3) match Table~\ref{tab:retrieval_multirepo}).

\begin{table}[!t]
\caption{Architectural Component-Wise Retrieval Ablations}
\label{tab:ablations}
\centering
\setlength{\tabcolsep}{3pt}
\resizebox{\columnwidth}{!}{%
\begin{tabular}{lcccc}
\toprule
\textbf{Component Variant} & \textbf{R@1} & \textbf{R@5} & \textbf{MRR} & \textbf{95\% CI (R@5)} \\
\midrule
\multicolumn{5}{l}{\textit{Django Codebase ($N=25$)}} \\
(1) Pure Lexical (TF-IDF) & 0.2000 & 0.6000 & 0.3667 & [0.407, 0.766] \\
(2) Pure AST (Tree-sitter) & 0.3200 & 0.6800 & 0.4447 & [0.484, 0.828] \\
(3) Static Hybrid RRF & \textbf{0.4400} & \textbf{0.8000} & \textbf{0.5773} & [0.609, 0.911] \\
(4) Unfiltered Episodic Transfer & 0.2000 & 0.3600 & 0.2460 & [0.203, 0.555] \\
\midrule
\multicolumn{5}{l}{\textit{Sphinx Codebase ($N=25$)}} \\
(1) Pure Lexical (TF-IDF) & 0.2400 & 0.4000 & 0.3033 & [0.234, 0.593] \\
(2) Pure AST (Tree-sitter) & 0.0000 & 0.1600 & 0.0600 & [0.064, 0.347] \\
(3) Static Hybrid RRF & 0.1600 & \textbf{0.4400} & 0.2447 & [0.267, 0.629] \\
(4) Unfiltered Episodic Transfer & 0.0000 & 0.1600 & 0.0513 & [0.064, 0.347] \\
\bottomrule
\end{tabular}%
}
\end{table}

Static Hybrid RRF achieves the best localization on Django (+57.4\% MRR over pure lexical). In contrast, accumulating episodic memory across arbitrary heterogeneous tasks without domain clustering causes \textbf{severe negative transfer} (MRR drops to 0.2460 on Django and 0.0513 on Sphinx). The mechanism is lexical: issue descriptions from unrelated subsystems share generic terms (e.g., \texttt{validation}, \texttt{field}, \texttt{model}), so past episodes touching unrelated files receive positive weight from surface overlap and outrank the true target in RRF fusion (see Case~4 for a worked analysis). Episodic transfer should therefore be scoped to recurrent subsystem clusters. We have not implemented or evaluated an automatic clustering gate; the scoping in Section~\ref{sec:clusters} was done manually.

\subsection{End-to-End Agent Task Execution ($N=3$, Answering RQ2)}
Table~\ref{tab:execution} summarizes full multi-turn agent runs inside the SWE-bench evaluation harness under matched model (\texttt{gemini-3.5-flash-lite}), turn budget (50 steps), and execution conditions.

\begin{table}[!t]
\caption{Master Execution Results Table ($N=3$ Sequential Tasks)}
\label{tab:execution}
\centering
\setlength{\tabcolsep}{3.5pt}
\resizebox{\columnwidth}{!}{%
\begin{tabular}{lcccc}
\toprule
\textbf{Condition} & \textbf{Resolved} & \textbf{Mean Turns} & \textbf{Mean Latency} & \textbf{Mean Cost} \\
\midrule
A (Stateless Baseline) & 1/3 (33.3\%) & 46.3 & 357.5s & \$0.0692 \\
B (Snapshot RAG) & 0/3 (0.0\%) & 40.7 & 264.5s & \$0.0757 \\
C (Naive Combo) & 1/3 (33.3\%) & 47.0 & 202.0s & \$0.0762 \\
\textbf{D (Proposed System)} & \textbf{1/3 (33.3\%)} & \textbf{48.7} & \textbf{207.3s} & \textbf{\$0.0679} \\
\bottomrule
\end{tabular}%
}
\end{table}

All resolving conditions resolved the identical task (\texttt{django-11848}). Despite placing target files at Rank 1 in Tasks 1 and 2, the proposed system did not resolve more tasks than the stateless baseline. \textbf{Answer to RQ2: no improvement in resolution was demonstrated.} Fault localization appears necessary but not sufficient for automated repair; with these sample sizes we cannot attribute the bottleneck to a specific cause, though trajectory inspection (Cases 2 and 3) points to multi-turn code synthesis and edit commitment.

\subsection{Cross-Task Episodic Transfer Across Recurrent Subsystems (Answering RQ5)}
\label{sec:clusters}
To test whether episodic memory transfers when scoped to related issues within the same subsystem, we evaluated two recurrent clusters (two task pairs in total):
\begin{enumerate}
    \item \textbf{Cluster 1 (Backend ORM)}: \texttt{django/db/models/base.py} (\texttt{django-12209} $\to$ \texttt{django-12273}).
    \item \textbf{Cluster 2 (Frontend Widgets)}: \texttt{django/forms/widgets.py} (\texttt{django-12193} $\to$ \texttt{django-12276}).
\end{enumerate}

Table~\ref{tab:clusters} reports the comparison between stateless Baseline A and persistent Condition D across both clusters.

\begin{table}[!t]
\caption{Cross-Subsystem Recurrent Cluster Evaluation}
\label{tab:clusters}
\centering
\setlength{\tabcolsep}{2.5pt}
\resizebox{\columnwidth}{!}{%
\begin{tabular}{lcccc}
\toprule
\textbf{Subsystem \& Task} & \textbf{Condition} & \textbf{Searches} & \textbf{Latency} & \textbf{Resolved} \\
\midrule
\multicolumn{5}{l}{\textit{Cluster 1: Backend Database ORM (\texttt{models/base.py})}} \\
Task 1: \texttt{django-12209} & A (Baseline) & 13 & 368.1s & False \\
Task 1: \texttt{django-12209} & D (Proposed) & 10 & 288.7s & False \\
Task 2: \texttt{django-12273} & A (Baseline) & 15 & 352.4s & False \\
Task 2: \texttt{django-12273} & \textbf{D (Proposed)} & \textbf{5 (-66.7\%)} & \textbf{262.1s} & False \\
\midrule
\multicolumn{5}{l}{\textit{Cluster 2: Frontend HTML Widgets (\texttt{forms/widgets.py})}} \\
Task 1: \texttt{django-12193} & A (Baseline) & 5 & 119.0s & \textbf{True} \\
Task 1: \texttt{django-12193} & \textbf{D (Proposed)} & \textbf{0 (-100\%)} & 246.4s & False \\
Task 2: \texttt{django-12276} & A (Baseline) & 4 & 361.6s & \textbf{True} \\
Task 2: \texttt{django-12276} & \textbf{D (Proposed)} & \textbf{4} & \textbf{283.6s (-21.6\%)} & \textbf{True} \\
\bottomrule
\end{tabular}%
}
\end{table}

In Cluster 1, persistent memory reduced exploratory search commands from 15 to 5 (-66.7\%) on Task 2, accelerating test reproduction to Turn 10, but neither condition resolved either task. In Cluster 2, memory routed directly to \texttt{widgets.py} at Rank 1, eliminating bash searches on Task 1 ($5 \to 0$), yet Condition D failed that task (Case 3). On Task 2 (\texttt{django-12276}), both conditions resolved the issue; Condition D did so in 41 turns and 283.6s, which is 78.0s (21.6\%) faster and 19.9\% lower in API cost than Baseline A. Across the four cluster tasks, Condition D resolved 1/4 and Baseline A resolved 2/4. Combined with the general sequence (Table~\ref{tab:execution}), this gives 2/7 for Condition D versus 3/7 for Baseline A over $N=7$ end-to-end tasks. Search reduction was therefore observed, but it did not translate into higher resolution, and the timing result rests on a single task.

\subsection{Incremental AST Caching Amortization (Answering RQ3)}
Table~\ref{tab:caching} tracks empirical cache performance across full parsing, warm lookup, and single-file drift.

\begin{table}[!t]
\caption{Incremental AST Caching and Invalidation Benchmarks}
\label{tab:caching}
\centering
\setlength{\tabcolsep}{3.5pt}
\resizebox{\columnwidth}{!}{%
\begin{tabular}{lccc}
\toprule
\textbf{Indexing State} & \textbf{Execution Time} & \textbf{Files Parsed} & \textbf{Cache Hits} \\
\midrule
Cold Baseline Parse & 2.076s & 804 / 804 & 0 / 804 (0.0\%) \\
Warm In-Memory Cache & \textbf{0.205s (10.13$\times$ speedup)} & 0 / 804 & 804 / 804 (\textbf{100.0\%}) \\
Partial Drift (1 File Modified) & 0.212s & 1 / 804 & 803 / 804 (99.9\%) \\
\bottomrule
\end{tabular}%
}
\end{table}

Warm cache lookups bypass redundant syntax tree traversal via $O(1)$ memory lookups backed by an $O(M)$ disk hash scan, delivering a \textbf{$10.13\times$ speedup}. When on-disk files drift, incremental invalidation re-parses only the affected file with negligible 7ms overhead.

\subsection{Systematic Repository Drift and Staleness Sensitivity (Answering RQ4)}
We evaluated memory behavior across four systematic repository drift scenarios:
\begin{enumerate}
    \item \textbf{Changed Function Signatures}: Target method signature modified with new parameters. Condition D detects the SHA-256 mismatch, re-indexes the AST incrementally, and extracts updated parameters; Condition C serves stale signatures.
    \item \textbf{Renamed / Moved Source Files}: Target file renamed. Condition D flags the missing file as deleted and purges invalid paths; Condition C continues to retrieve the deleted path.
    \item \textbf{Deleted Symbols}: Deprecated function removed. Condition D purges the symbol from AST cache, preventing API hallucination.
    \item \textbf{Conflicting Historical Memories}: Earlier episode touched a drifted file while another touched an untouched file. Condition D selectively attenuates the drifted episode while preserving full reinforcement for fresh episodes.
\end{enumerate}

\begin{table}[!t]
\caption{Staleness Attenuation Factor ($\lambda$) Sensitivity Analysis}
\label{tab:staleness_sweep}
\centering
\setlength{\tabcolsep}{2.5pt}
\resizebox{\columnwidth}{!}{%
\begin{tabular}{lccc>{\raggedright\arraybackslash}p{3.4cm}}
\toprule
$\boldsymbol{\lambda}$ & \textbf{MRR} & \textbf{Hit@1} & \textbf{Hit@5} & \textbf{Observation} \\
\midrule
$\lambda = 0.00$ & 0.6667 & 0.6667 & 0.6667 & Hard eviction; discards valid hints \\
$\lambda = 0.10$ & 1.0000 & 1.0000 & 1.0000 & Saturated on this scenario set \\
$\boldsymbol{\lambda = 0.20}$ & \textbf{1.0000} & \textbf{1.0000} & \textbf{1.0000} & \textbf{Default; conservative choice} \\
$\lambda = 0.30$ & 1.0000 & 1.0000 & 1.0000 & Saturated on this scenario set \\
$\lambda = 0.50$ & 1.0000 & 1.0000 & 1.0000 & Saturated on this scenario set \\
$\lambda = 0.80$ & 1.0000 & 1.0000 & 1.0000 & Saturated; pollution risk not observed \\
$\lambda = 1.00$ & 1.0000 & 1.0000 & 1.0000 & No penalty; not separated from above \\
\bottomrule
\end{tabular}%
}
\end{table}

Table~\ref{tab:staleness_sweep} shows that hard eviction ($\lambda = 0.00$) causes premature amnesia, dropping MRR from 1.0000 to 0.6667 by discarding useful module hints. Every $\lambda \ge 0.10$ achieves identical, saturated performance on our four scenarios, so this sweep \emph{cannot} distinguish between them. In particular, it does not demonstrate that weak penalties ($\lambda \ge 0.8$) pollute prompts; that remains a plausible but untested risk. We select $\lambda = 0.20$ as a conservative default that avoids hard eviction while strongly down-weighting stale episodes, and leave a larger, more discriminating scenario set to future work.

\section{Qualitative Case Studies, Post-Mortem \& Negative Transfer}

\subsection{Case 1: Resolved Bug in \texttt{django-11848}}
RFC 7231 Section 7.1.1.1 mandates that HTTP date parsers interpreting 2-digit years must apply a sliding window assigning dates over 50 years in the future to the previous century. Under Condition D, the agent received repository context routing directly to \texttt{django/utils/http.py} (\texttt{parse\_http\_date}) at Rank 2. The agent navigated directly to the target at turn 2, implemented dynamic century windowing, verified the fix against unit tests, and successfully passed all SWE-bench harness assertions.

\subsection{Case 2: Downstream Synthesis Failure in \texttt{django-11815}}
In \texttt{django-11815}, the migration serializer incorrectly wrote \texttt{Enum.name} instead of \texttt{Enum.value}. Memory accurately localized \texttt{django/db/migrations/serializer.py} and class \texttt{EnumSerializer} at Rank 1. However, the agent's multi-turn regex edits introduced syntax errors and failed to converge before turn exhaustion, highlighting the \textbf{Localization--Resolution Disconnect}.

\subsection{Case 3: Post-Mortem of Frontend Failure in \texttt{django-12193}}
In the frontend widgets cluster (Table~\ref{tab:clusters}), stateless Baseline A resolved \texttt{django-12193} in 44 turns, whereas Condition D failed with an empty patch after 37 turns despite Rank-1 localization ($0$ search commands). 

Trajectory inspection suggests an interaction between memory context and prompt budget pacing: Baseline A spent 5 exploratory bash commands searching files, which framed the task around a single localized diff in \texttt{widgets.py}. In Condition D, rich AST context provided all class hierarchies (\texttt{Widget}, \texttt{TextInput}, \texttt{CheckboxInput}). The agent immediately entered an intensive reproduction testing phase at turn 10, exploring interactions between \texttt{Widget.build\_attrs} and \texttt{CheckboxInput.get\_context}. However, as the turn count approached the budget limit, the agent over-deliberated without committing the edit, and budget-pacing prompts induced an empty patch submission. This single trajectory suggests that \textbf{dense structural memory can induce over-deliberation in reproduction testing}; it is a hypothesis from one case, which execution scaffolding enforcing strict editing checkpoints could test.

\subsection{Case 4: Negative Transfer in Unfiltered Episodic Streams}
When episodic memory was updated continuously across the heterogeneous 25-task Django sequence without subsystem clustering (Table~\ref{tab:ablations}), localization MRR collapsed from 0.5773 to 0.2460. Lexical similarity between unrelated issue descriptions (e.g., generic terms like \texttt{validation}, \texttt{field}, \texttt{model}) caused past episodes from unrelated subsystems to outrank true target files in RRF fusion. This is empirical evidence of \textbf{episodic negative transfer}: episodic memory is beneficial only when scoped to recurrent subsystems, and would need to be gated by domain-clustering heuristics in production. Under task-specific chronological checkouts the same effect appears across all 50 tasks: unclustered accumulation lowers combined MRR from 0.5213 to 0.1860.

\section{Discussion: When Memory Helps and When It Hurts}
Our results show that persistent repository memory is not uniformly beneficial. We identify four distinct ways in which memory was neutral or harmful.

\textbf{Unscoped episodic memory (negative transfer).} Episodes retrieved by lexical overlap import irrelevant files into ranking. The damage is large (MRR 0.5773$\to$0.2460 on Django, 0.5213$\to$0.1860 combined) and arises because generic vocabulary is shared across subsystems. Memory must be scoped, and an automatic scoping mechanism remains unimplemented.

\textbf{Context that is too rich (over-deliberation).} In \texttt{django-12193}, accurate localization plus dense class-hierarchy context coincided with an empty patch, while the less-informed baseline resolved the task. This is a single observation and its cause is untested, but it indicates that better context can change agent behavior in ways retrieval metrics do not capture, consistent with the Localization--Resolution Disconnect.

\textbf{Stale memory.} SHA-256 tracking invalidated obsolete signatures, renamed files, and deleted symbols in all four drift scenarios, whereas the naive combination served stale information. However, invalidation is file-granular, so an unrelated edit anywhere in a stored file attenuates the entire episode, which may discard valid hints. Hard eviction ($\lambda = 0$) demonstrably harms retrieval, but our scenario set is too small to separate $\lambda$ values above 0.1.

\textbf{Retrieval-modality mismatch.} AST symbol matching depends on issue text naming code identifiers; on Sphinx it was nearly useless alone (MRR 0.0600), and fusion did not beat lexical retrieval on MRR under the snapshot protocol. Memory built on code structure inherits this dependence.

Taken together, memory components should be treated as inputs whose benefit must be measured against the end task, not only against retrieval accuracy.

\section{Threats to Validity}
\textbf{Temporal Validity}: The primary retrieval results use each task's exact \texttt{base\_commit}, avoiding future-information leakage. The ablation and episodic-transfer experiments use the single pinned snapshot (\texttt{b1d6b35e}) and may therefore differ from the chronological results.

\textbf{End-to-End Sample Size}: End-to-end multi-turn evaluation covers $N=7$ task executions per condition ($N=3$ general + $N=4$ cluster tasks) with a single run each. This is too small to establish or rule out differences in resolution rate; localization claims rest on $N=50$ tasks, resolution observations on $N=7$.

\textbf{Repository Diversity}: Only two repositories (Django and Sphinx) were evaluated, and both are Python projects.

\textbf{Model Capability Ceiling}: Experiments used \texttt{gemini-3.5-flash-lite}; frontier models with stronger synthesis capabilities may convert Rank-1 localization into higher resolution rates, or may behave differently under dense context.

\textbf{Baselines}: Comparisons are against lexical and pure-AST retrieval and a stateless agent; we do not compare against other published repository-memory or repository-RAG systems (Table~\ref{tab:comparison} is qualitative).

\section{Reproducibility Evidence}
All experimental artifacts are packaged in the open-source reproducibility archive \texttt{reproducibility\_package.tar.gz} (0.81\,MB), containing: (1) full executable source code (\texttt{memory\_pilot/}); (2) task lists and SWE-bench configurations (\texttt{inventory.json}); (3) raw benchmark JSONs; (4) per-task logs, patches (\texttt{patch.diff}), and Docker evaluation reports; and (5) the standalone script \texttt{scripts/run\_statistical\_tests.py}, allowing exact reproduction of all manuscript statistics, confidence intervals, and Wilcoxon tests in a single command.

\section{Conclusion}
We evaluated a \textbf{Session-Persistent Repository Memory Architecture} for SWE agents. Fusing Tree-sitter AST symbol graphs with sparse lexical retrieval via RRF improves fault localization on Django (MRR 0.6027 vs.\ 0.3667; 80\% Top-5 hit rate) and on the combined 50-task benchmark (MRR 0.5213 vs.\ 0.4163; 70\% Top-5), with mixed results on Sphinx. SHA-256 caching delivers a $10.13\times$ indexing speedup, and SHA-256 staleness tracking handled all four drift scenarios tested. In two recurrent Django clusters, scoped episodic memory reduced exploratory searches, but we did not demonstrate improved task resolution (2/7 versus 3/7 for the stateless baseline), and unscoped episodic memory caused severe negative transfer. These results show that retrieval gains do not by themselves yield better coding performance. Establishing task-level benefits will require more repositories, larger end-to-end evaluations with stronger models, comparisons against existing memory systems, and an implemented mechanism for scoping episodic memory.

\bibliographystyle{IEEEtran}
\bibliography{references}

\end{document}
